\section{Introduction}
\label{sec:intro}

Suppose you have just read a public vulnerability fix for a C library you
depend on: a one-line patch that frees a buffer and then clears the
pointer to it, because the code path that frees it can be reached again
before the function returns. You trust the fix. The natural next question
is whether the same mistake, the same missing reset after the same kind
of release, sits somewhere else in the same codebase, in a function the
patch never touched.

Static analyzers already look for memory-safety bugs in C, and some are
free and widely used: Clang Static Analyzer, Cppcheck, and Coccinelle
semantic patches, among others. These tools ship with general rule sets
built from common bug patterns across many projects. They stop short of
one thing: a rule derived from a single project's own patch history, for
the exact ownership or reset convention that patch teaches, checked
everywhere that convention should hold in that project's current source.

We focus on a second problem that starts where general-purpose rule sets
stop. Given one historical fix, can a rule extracted from it, with no
function names baked in, be checked against a project's whole source
tree? We ask three things of such a rule: (i)~flag the function the fix
came from before the fix, and clear it after; (ii)~do the same for an
unrelated fix nobody built the rule from; and (iii)~say, for every site it
flags on the project's current source, whether that site is conforming, a
real violation, or explainable, with a witness a human can check by hand.

CVE-2022-4450~\cite{nvd2022cve4450} is a real case. \texttt{PEM\_read\_bio\_ex} in OpenSSL's
\texttt{crypto/pem/pem\_lib.c} frees two output buffers on one failure
path (Figure~\ref{fig:motivating}) without clearing the caller-visible
pointers to them; a caller that frees these outputs again after failure
can double-free. The fix, commit \texttt{63bcf189be}, adds exactly two
assignments.

Closest to our rule-as-data approach is VulPA~\cite{cao2025vulpa}, which
also derives a checkable rule from a known bug, but compiles it to an
automaton checked by an interprocedural analysis with aliasing, on Java.
Ours is intraprocedural, with no alias analysis (Section~\ref{sec:related}).

We present \tool, a checker that reads a lifecycle rule from a data file
and applies it to real C source through a parser and a per-function
control-flow graph built with \texttt{tree-sitter}, searches for a path
from each release of an output parameter to function exit that never
resets it, and classifies each site with a witness path
(Section~\ref{sec:approach}). This paper implements two such rules against the same parser and CFG
builder: Protocol~B, for the output-pointer reset convention behind
CVE-2022-4450, and Protocol~A, for the detach-before-release BIO-chain
convention behind CVE-2023-0215~\cite{nvd2023cve0215} (\rqref{5}, Section~\ref{sec:rq5}).

On OpenSSL, \tool flags \texttt{PEM\_read\_bio\_ex}'s two output
parameters before the fix and clears both after it. It also reproduces an
unrelated non-CVE fix the rule was never shown, and flags one site a
later OpenSSL commit went on to fix independently
(Section~\ref{sec:eval}). On the current source tree it flags 10 sites.
Manual triage finds 2 potential recurrences whose failure branch is not
reachable in normal builds, 8 harmless deviations, and 0
confirmed new vulnerabilities. The lesson we did not expect going in:
none of three general-purpose open-source static analyzers, run with
their stock configuration, distinguishes any of the three CVEs we
checked from its own fix (Section~\ref{sec:rq2}). A project-specific
rule, not a sharper general one, is what closed the gap here. Protocol~A
also distinguishes its own motivating fix, but its whole-tree run finds
0 confirmed recurrences out of 11 flagged sites, all false positives from
the rule's missing notion of object ownership: the rule-as-data
architecture generalizes across invariant shapes even where a given
rule's detection does not generalize across a whole codebase
(Section~\ref{sec:rq5}).

\paragraph{Contributions.} (1)~A rule-as-data lifecycle invariant for
output-pointer reset after release, independent of any function name
(Section~\ref{sec:problem}); (2)~\tool, which checks it on a real
\texttt{tree-sitter} control-flow graph (Section~\ref{sec:approach});
(3)~an evaluation on six pinned OpenSSL snapshots and a four-tool
comparison on three historical CVEs (Section~\ref{sec:eval}); (4)~a second,
differently-shaped rule, Protocol~A, showing where single-object
generalization breaks down (Section~\ref{sec:rq5}); and (5)~the released
tool and data (Section~\ref{sec:data}).

\begin{figure}[t]
\begin{lstlisting}[numbers=none,escapechar=@]
headerlen = BIO_get_mem_data(headerB, NULL);
*header = pem_malloc(headerlen + 1, flags);
*data = pem_malloc(len, flags);
if (*header == NULL || *data == NULL)
    goto out_free;
...
out_free: @\textcolor{red}{// *header, *data freed, never reset}@
    pem_free(*header, flags, 0);
    pem_free(*data, flags, 0);
end:
    ...
    return ret; @\textcolor{red}{// caller sees stale *header, *data}@
\end{lstlisting}
\caption{\texttt{PEM\_read\_bio\_ex} in OpenSSL \texttt{8e257b86e5}
  (\texttt{crypto/pem/pem\_lib.c}, trimmed). The \texttt{out\_free} path
  frees \texttt{*header} and \texttt{*data} but never resets either
  pointer before \texttt{return}. The fix, commit \texttt{63bcf189be},
  adds \texttt{*header = NULL;} and \texttt{*data = NULL;} to this path.}
\label{fig:motivating}
\Description{Code excerpt from OpenSSL showing two output pointers freed on an error path without being reset to null before the function returns.}
\end{figure}
